\chapter{Allergy Blood Testing (Specific IgE)}

\section*{What Is This Test?}
Allergy blood testing measures allergen-specific immunoglobulin E (sIgE) antibodies in serum to identify sensitization to specific allergens. IgE is the class of antibody that mediates type I (immediate) hypersensitivity reactions, also known as allergic reactions. When an allergen cross-links two IgE molecules bound to the high-affinity IgE receptor (FceRI) on mast cells and basophils, it triggers degranulation and the release of histamine, leukotrienes, prostaglandins, and other inflammatory mediators, producing the symptoms of allergic disease: urticaria, angioedema, rhinoconjunctivitis, bronchospasm, and anaphylaxis. Specific IgE blood testing (formerly known as RAST, radioallergosorbent test) uses modern immunoassay platforms to quantify IgE directed against individual allergens including pollens, animal danders, dust mites, molds, insect venoms, foods, drugs, and occupational allergens. The test provides an alternative to skin prick testing and is particularly valuable when skin testing is contraindicated (severe eczema, dermatographism, antihistamine use) or when the risk of systemic allergic reaction from skin testing is high.

\section*{Alternative Names}
Allergen-Specific IgE; sIgE; IgE RAST; ImmunoCAP Specific IgE; Allergy Blood Test; In Vitro Allergy Testing; Antigen-Specific IgE; Allergen-Specific Immunoglobulin E

\section*{Principle}
Allergen-specific IgE testing is performed using the ImmunoCAP (Thermo Fisher Scientific/Phadia) system or other commercial immunoassay platforms (Immulite, Hycor). In the ImmunoCAP system, allergen of interest is covalently coupled to a cellulose solid phase matrix encased in a plastic cap. Patient serum is incubated with the allergen-coated matrix, allowing allergen-specific IgE antibodies to bind. After washing, enzyme-conjugated (beta-galactosidase) anti-human IgE detection antibody is added. Following a final wash, a fluorogenic substrate (4-methylumbelliferyl-beta-D-galactoside) is added and the resulting fluorescence is measured. The fluorescence signal is directly proportional to the allergen-specific IgE concentration and is calibrated against a standard curve referenced to the WHO International Standard for Human IgE (WHO 75/502). Results are reported quantitatively in kU$_\text{A}$/L (kilo allergy units per liter). The lower limit of quantitation is typically 0.1 kU$_\text{A}$/L, with results $\geq$0.35 kU$_\text{A}$/L considered positive. Total IgE is often measured simultaneously. Multiplex platforms (e.g., ImmunoCAP ISAC, ALEX) can measure specific IgE to over 100 allergen components simultaneously using microarray technology.

\section*{History}
The discovery of IgE as the reaginic antibody mediating allergic reactions was made independently by Kimishige and Teruko Ishizaka in the United States (1966) \cite{ishizaka1966physicochemical} and by S.G.O. Johansson and Hans Bennich in Sweden (1967). The radioallergosorbent test (RAST), the first in vitro assay for allergen-specific IgE, was developed by Wide, Bennich, and Johansson in 1967, using radiolabeled anti-IgE antibodies for detection \cite{wide1967diagnosis}. The transition from RAST (radioactivity-based) to enzyme- and fluorescence-based detection (ImmunoCAP, introduced in 1989) improved sensitivity, precision, and safety. The development of recombinant allergen components in the 1990s--2000s enabled component-resolved diagnostics (CRD), which differentiates between genuine sensitization (to species-specific allergen components) and cross-reactivity (to panallergens such as profilins or cross-reactive carbohydrate determinants). The WHO/IUIS Allergen Nomenclature Sub-Committee maintains the official nomenclature for allergen sources and components.

\section*{Why This Test Is Ordered}
\begin{itemize}
  \item Diagnosis of IgE-mediated allergic disease in patients with allergic rhinitis, allergic conjunctivitis, allergic asthma, atopic dermatitis, urticaria, angioedema, food allergy, insect venom allergy, drug allergy, latex allergy, or occupational allergy
  \item Identification of specific allergens triggering symptoms when the history suggests multiple possible triggers or when the trigger is unclear
  \item Evaluation of food allergy when skin prick testing is contraindicated (severe eczema, dermatographism, inability to discontinue antihistamines) or when the risk of systemic reaction from skin testing is high
  \item Component-resolved diagnosis (CRD) to distinguish true food allergy from cross-reactivity (e.g., Ara h 2 specific IgE for peanut allergy vs. profilin or CCD cross-reactivity) and to assess risk of severe reactions
  \item Monitoring of food allergy resolution, particularly in children with milk, egg, soy, and wheat allergy, where falling specific IgE levels may predict development of tolerance
  \item Evaluation of patients with recurrent idiopathic anaphylaxis in whom no trigger has been identified by history
  \item Confirmation of venom allergy (honeybee, yellow jacket, paper wasp, fire ant) before initiation of venom immunotherapy
  \item Evaluation of suspected latex allergy in healthcare workers, spina bifida patients, and individuals with multiple surgeries
  \item Quantification of allergen-specific IgE to guide allergen immunotherapy decisions and to monitor the immunologic response to treatment
\end{itemize}

\section*{Is This a Primary or Secondary Test}
This is a primary diagnostic test for IgE-mediated allergic disease, used as an alternative or complement to skin prick testing. In clinical practice, skin prick testing is often performed first because it provides immediate results, is more sensitive for some allergens, and covers a broader panel of allergens at a lower cost. However, specific IgE blood testing is equally valid and is preferred in specific clinical situations (severe eczema, antihistamine use, anaphylaxis risk).

\section*{How This Test Is Performed}
For most patients, a health care professional will take a blood sample from a vein in your arm using a small needle. After the needle is inserted, a small amount of blood will be collected into a test tube or vial. You may feel a little sting when the needle goes in or out. This usually takes less than five minutes.

\section*{What Preparation Is Needed}
No special preparation is required. Unlike skin prick testing, antihistamines, tricyclic antidepressants, and other medications that interfere with skin testing do NOT affect allergen-specific IgE blood test results and do not need to be discontinued. Fasting is not necessary. Patients should inform their healthcare provider about any recent systemic allergic reactions, as IgE levels may be transiently consumed during a severe allergic episode and produce a false-negative result. It is generally recommended to wait 2--4 weeks after a severe anaphylactic episode before performing specific IgE testing.

\section*{Contraindications and Precautions}
No absolute contraindications to venipuncture. Specimen should be collected in a serum separator tube (SST, gold-top) or red-top tube. Specific IgE blood test results must always be interpreted in the context of clinical history. A positive specific IgE indicates sensitization (immune response to the allergen), NOT necessarily clinical allergy. The positive predictive value of a specific IgE test varies by allergen: for some allergens (peanut, egg, milk, shrimp), a positive specific IgE at a certain threshold has a high predictive value for clinical allergy, while for other allergens (some pollens, grasses) a positive result more commonly represents asymptomatic sensitization. Cross-reactivity between homologous allergens from related species (e.g., tree pollens, grass pollens, legumes, shellfish) can lead to multiple positive results without corresponding clinical allergy to each cross-reactive allergen. Component-resolved diagnostics improve diagnostic specificity by distinguishing genuine sensitization from cross-reactivity. Very high total IgE levels (>1,000--2,000 kU/L) can cause non-specific binding and false-positive specific IgE results. Patients with immunodeficiency or immunosuppression may have false-negative results due to impaired IgE production. Eosinophilic gastrointestinal disorders may present with negative specific IgE despite clinical food-triggered symptoms, as these are often cell-mediated (non-IgE-mediated) reactions.

\section*{Risks}
Minimal risk. Slight pain or bruising at the venipuncture site. Rare: hematoma, infection, vasovagal syncope.

\section*{What Happens After the Test}
Results are typically available within 3 to 7 days (longer for component testing). Results are reported quantitatively in kU$_\text{A}$/L for each allergen tested. The magnitude of specific IgE elevation correlates with the probability of clinical reactivity, but there is no absolute threshold above which clinical allergy can be assumed. Results should be interpreted by a physician with expertise in allergic disease (allergist-immunologist), who will integrate the specific IgE results with the clinical history, physical examination, and (when indicated) the results of skin prick testing and oral food challenge.

\section*{Understanding Results}

\subsection*{Below Normal (Negative or Undetectable)}
\begin{itemize}
  \item Negative specific IgE (<0.35 kU$_\text{A}$/L) to a clinically suspected allergen: Does not exclude clinical allergy with 100\% certainty. False-negative results occur due to: IgE consumption during a recent allergic reaction (test 2--4 weeks after reaction), sensitization to an allergen not included in the test panel, sensitivity to a minor allergen component not adequately represented in the extract, or local allergic disease (allergic rhinitis confined to the nasal mucosa without systemic IgE sensitization)
  \item Localized allergic rhinitis (entopy): Patients with typical allergic rhinitis symptoms and negative skin prick testing and negative specific IgE may have IgE production confined to the nasal mucosa without systemic detection. Nasal allergen provocation test confirms the diagnosis
  \item Food protein-induced enterocolitis syndrome (FPIES) and food protein-induced allergic proctocolitis (FPIAP): Non-IgE-mediated food allergies presenting with delayed gastrointestinal symptoms (vomiting, diarrhea, failure to thrive). Specific IgE is typically negative. Diagnosis is based on clinical history and resolution with dietary elimination
  \item Low or undetectable specific IgE to milk, egg, soy, or wheat in a young child with a previous food allergy history may indicate that the allergy has been outgrown and that an oral food challenge may be appropriate to confirm resolution
  \item Outgrowing of sensitization: Falling specific IgE levels over time, particularly in children, suggest developing tolerance. A reduction of 50\% or more over 12--24 months is a favorable prognostic sign
\end{itemize}

\subsection*{Above Normal (Positive)}
\begin{itemize}
  \item Sensitization (positive specific IgE) does NOT equal clinical allergy. A positive specific IgE result must be correlated with clinical history of symptoms on exposure to that allergen. Asymptomatic sensitization is common, particularly for aeroallergens (pollen, dust mite, animal dander) and some food allergens
  \item Food allergy: Specific IgE levels correlate with the probability of clinical food allergy \cite{sampson2014food}. 95\% positive predictive value (PPV) thresholds (ImmunoCAP) have been established for certain foods in children: egg >7 kU$_\text{A}$/L (infants <2 years: >2 kU$_\text{A}$/L), milk >15 kU$_\text{A}$/L (infants <1 year: >5 kU$_\text{A}$/L), peanut >14 kU$_\text{A}$/L, fish >20 kU$_\text{A}$/L, soy >30 kU$_\text{A}$/L, wheat >26 kU$_\text{A}$/L. Results above these thresholds predict a >95\% likelihood of clinical reactivity on oral food challenge. A result below the 95\% PPV threshold does not rule out clinical allergy
  \item Component-resolved diagnosis (CRD): Specific IgE to allergen components refines the diagnosis. For peanut allergy: elevated Ara h 2 specific IgE is highly predictive of systemic (anaphylactic) peanut allergy, whereas elevated Ara h 8 (a Bet v 1 homologue) specific IgE typically indicates mild oral allergy symptoms due to birch pollen cross-reactivity without risk of anaphylaxis. For egg allergy: Gal d 1 (ovomucoid) specific IgE indicates persistent allergy with heat-stable allergen; isolated Gal d 2 (ovalbumin) indicates a heat-labile allergen and tolerance of extensively baked egg. For milk allergy: Bos d 8 (casein) specific IgE indicates persistent allergy with heat-stable allergen
  \item Aeroallergen sensitization: Positive specific IgE to pollens (birch, grass, ragweed, mugwort), dust mites (Dermatophagoides pteronyssinus, D. farinae), animal danders (cat Fel d 1, dog Can f 1), and molds (Alternaria, Aspergillus, Cladosporium) supports the diagnosis of seasonal or perennial allergic rhinitis, allergic conjunctivitis, and allergic asthma
  \item Venom allergy: Positive specific IgE to honeybee venom (Api m 1), yellow jacket venom (Ves v 1 and Ves v 5), and paper wasp venom (Pol d 5) supports the diagnosis of systemic venom allergy. Component-resolved diagnosis (CRD) helps distinguish true venom sensitization from cross-reactivity due to cross-reactive carbohydrate determinants (CCDs)
  \item Drug allergy: Specific IgE testing is available for a limited number of drugs including penicillin major and minor determinants (penicilloyl G, penicilloyl V, penicilloate, penilloate), amoxicillin, and beta-lactam antibiotics. Negative specific IgE does not exclude drug allergy, and drug challenge may be required
  \item Latex allergy: Elevated specific IgE to latex (Hevea brasiliensis) in a patient with clinical reactions on latex exposure (contact urticaria, angioedema, asthma, anaphylaxis). Recombinant latex allergen components (Hev b 1, Hev b 3, Hev b 5, Hev b 6.01, Hev b 8) improve diagnostic specificity and help distinguish occupational latex allergy from pollen-food cross-reactivity
  \item Total IgE elevation with specific IgE: Elevated total IgE (>100 IU/mL) supports an atopic constitution (atopic dermatitis, allergic asthma, allergic rhinitis). Markedly elevated total IgE (>1,000--10,000 IU/mL) may be seen in severe atopic dermatitis, hyper-IgE syndrome (STAT3 deficiency with eczema and recurrent staphylococcal abscesses), allergic bronchopulmonary aspergillosis (ABPA), and parasitic infections
\end{itemize}

\section*{Normal Results}
\begin{itemize}
  \item \textbf{Allergen-Specific IgE:}
  \begin{itemize}
    \item Negative: <0.35 kU$_\text{A}$/L (class 0)
    \item Equivocal/Low positive: 0.35--0.70 kU$_\text{A}$/L (class 1)
    \item Moderate positive: 0.71--3.50 kU$_\text{A}$/L (class 2)
    \item High positive: 3.51--17.5 kU$_\text{A}$/L (class 3)
    \item Very high positive: 17.6--50 kU$_\text{A}$/L (class 4)
    \item Extremely high positive: >50 kU$_\text{A}$/L (class 5 and 6)
  \end{itemize}
  \item \textbf{Total IgE (age-dependent reference ranges):}
  \begin{itemize}
    \item Infants (0--6 months): <13 IU/mL
    \item Children (6--12 months): <34 IU/mL
    \item Children (1--3 years): <53 IU/mL
    \item Children (3--5 years): <100 IU/mL
    \item Children (5--9 years): <127 IU/mL
    \item Children (9--14 years): <200 IU/mL
    \item Adults: <100 IU/mL
  \end{itemize}
  \item Total IgE levels are not specific for any allergen and must be interpreted in conjunction with allergen-specific IgE testing
  \item The clinical significance of a positive specific IgE result depends on both the magnitude of the result and the specific allergen, as predictive values vary widely
\end{itemize}

\section*{Follow-up Tests}
\begin{itemize}
  \item \textbf{Skin prick testing (SPT):} The primary alternative method for detecting allergen sensitization. Intradermal introduction of allergen extracts into the skin produces a wheal-and-flare response in sensitized individuals within 15--20 minutes. More rapid results than serum specific IgE, with comparable sensitivity. Requires discontinuation of antihistamines (5--7 days) and is contraindicated in patients with severe eczema or dermatographism
  \item \textbf{Component-resolved diagnosis (ImmunoCAP ISAC or ALEX):} Microarray-based testing for specific IgE to over 100 individual allergen components (molecules) simultaneously. Distinguishes genuine sensitization from cross-reactivity, assesses risk of severe reactions (e.g., Ara h 2 for peanut anaphylaxis, Cor a 8 for hazelnut systemic reactions), and identifies cross-reactive patterns to guide avoidance and immunotherapy
  \item \textbf{Total IgE:} Provides baseline information about total IgE production and atopic status. Elevated total IgE supports the diagnosis of atopic disease and is also used in the diagnostic criteria for allergic bronchopulmonary aspergillosis (ABPA) and hyper-IgE syndrome
  \item \textbf{Oral food challenge (OFC):} The gold standard for diagnosing food allergy. Administered under medical supervision in a facility equipped to manage anaphylaxis. Used when the clinical history and test results are equivocal or when determining if a food allergy has been outgrown
  \item \textbf{Drug provocation test (DPT, drug challenge):} Graded administration of a suspected drug under medical supervision when specific IgE testing and skin testing are negative or unavailable but clinical suspicion remains
  \item \textbf{Pulmonary function tests (spirometry):} Including pre- and post-bronchodilator spirometry, to evaluate for allergic asthma in patients with aeroallergen sensitization and respiratory symptoms. Methacholine challenge if spirometry is normal but clinical suspicion remains
  \item \textbf{Basophil activation test (BAT):} A functional assay that measures basophil degranulation (CD63 and CD203c expression) after in vitro allergen stimulation. BAT can be useful in patients with discordant results between clinical history, skin testing, and serum specific IgE. Primarily available in research and specialized laboratories
  \item \textbf{Nasal allergen provocation test:} Used for diagnosis of local allergic rhinitis in patients with typical symptoms and negative skin testing and serum specific IgE
  \item \textbf{Serum tryptase:} Measured during an acute systemic allergic reaction (within 1--3 hours) to confirm mast cell activation. An elevated serum tryptase during an acute episode with a normal baseline level confirms anaphylaxis
\end{itemize}

\section*{References}
\bibliographystyle{unsrtnat}
\bibliography{references}

\clearpage
